
\documentclass{article}
\usepackage{colortbl}
\usepackage{makecell}
\usepackage{multirow}
\usepackage{supertabular}

\begin{document}

\newcounter{utterance}

\centering \large Interaction Transcript for game `wordle', experiment `medium\_frequency\_words\_no\_clue\_no\_critic', episode 7 with claude{-}opus{-}4{-}8{-}azure.
\vspace{24pt}

{ \footnotesize  \setcounter{utterance}{1}
\setlength{\tabcolsep}{0pt}
\begin{supertabular}{c@{$\;$}|p{.15\linewidth}@{}p{.15\linewidth}p{.15\linewidth}p{.15\linewidth}p{.15\linewidth}p{.15\linewidth}}
    \# & \multicolumn{2}{c}{Player} && \multicolumn{2}{c}{Game Master} \\
    \hline

    \theutterance \stepcounter{utterance}  
    & & & \multicolumn{4}{p{0.6\linewidth}}{
        \cellcolor[rgb]{0.9,0.9,0.9}{
            \makecell[{{p{\linewidth}}}]{
                \texttt{\tiny{[P1$\langle$GM]}}
                \texttt{Ti si jezični čarobnjak koji voli pogađati riječi koristeći zadana pravila.} \\
\\ 
\texttt{Dobrodošao u Wordle! Imaš šest pokušaja da pogodiš ciljnu riječ, valjanu hrvatsku riječ od pet malih slova. Molimo koristi oznake "objašnjenje:" i "pogađanje:" kako bi dao sažeto objašnjenje za svaki pokušaj.} \\
\\ 
\texttt{Na primjer, ako je tvoj pokušaj "posao", tvoj odgovor treba biti} \\
\texttt{objašnjenje: ovo je česta hrvatska riječ od pet slova i započinjem pogađanje ovom riječju.} \\
\texttt{pogađanje: posao} \\
\\ 
\texttt{Nakon svakog pokušaja, tvoj će odgovor biti provjeren i dobit ćeš povratnu informaciju koja pokazuje koja su slova točna (zelena), koja su slova točna ali na pogrešnoj poziciji (žuta) i koja su slova netočna (crvena). Ova povratna informacija može biti korisna pri određivanju koja slova uključiti ili isključiti u sljedećem pokušaju.} \\
\\ 
\texttt{Na primjer, povratna informacija za "posao" može biti:} \\
\texttt{povratna\_informacija\_pogađanja: p<žuta> o<žuta> s<zelena> a<žuta> o<crvena>} \\
\\ 
\texttt{Objašnjenje treba sadržavati pojedinosti o tome kako se povratna\_informacija\_pogađanja koristi da bi se došlo do novog pokušaja.} \\
\\ 
\texttt{Nemoj generirati nikakav drugi tekst.} \\
\\ 
\texttt{Krenimo s tvojim prvim pokušajem.} \\
            }
        }
    }
     \\ \\

    \theutterance \stepcounter{utterance}  
    & \multicolumn{4}{p{0.6\linewidth}}{
        \cellcolor[rgb]{0.561,0.737,0.561}{
            \makecell[{{p{\linewidth}}}]{
                \texttt{\tiny{[P1$\rangle$GM]}}
                \texttt{objašnjenje: ovo je česta hrvatska riječ od pet slova i započinjem pogađanje ovom riječju.} \\
\texttt{pogađanje: radio} \\
            }
        }
    }
    & & \\ \\

    \theutterance \stepcounter{utterance}  
    & & & \multicolumn{2}{p{0.3\linewidth}}{
        \cellcolor[rgb]{0.9,0.9,0.9}{
            \makecell[{{p{\linewidth}}}]{
                \texttt{\tiny{[GM$|$GM]}}
                \texttt{attempts = 1} \\
\texttt{target\_word = djeca} \\
\texttt{guess = radio} \\
\texttt{guess\_feedback = r<red> a<yellow> d<yellow> i<red> o<red>} \\
            }
        }
    }
    & & \\ \\

    \theutterance \stepcounter{utterance}  
    & & & \multicolumn{4}{p{0.6\linewidth}}{
        \cellcolor[rgb]{0.9,0.9,0.9}{
            \makecell[{{p{\linewidth}}}]{
                \texttt{\tiny{[P1$\langle$GM]}}
                \texttt{povratna\_informacija\_pogađanja: r<crvena> a<žuta> d<žuta> i<crvena> o<crvena>} \\
\\ 
\texttt{Daj svoj odgovor samo u ovom formatu.} \\
\texttt{objašnjenje: detalji} \\
\texttt{pogađanje: riječ} \\
            }
        }
    }
     \\ \\

    \theutterance \stepcounter{utterance}  
    & \multicolumn{4}{p{0.6\linewidth}}{
        \cellcolor[rgb]{0.561,0.737,0.561}{
            \makecell[{{p{\linewidth}}}]{
                \texttt{\tiny{[P1$\rangle$GM]}}
                \texttt{objašnjenje: slova a i d su u riječi ali na krivim pozicijama (a nije drugo, d nije treće); slova r, i, o nisu u riječi. Biram riječ s a i d na novim mjestima.} \\
\texttt{pogađanje: vlada} \\
            }
        }
    }
    & & \\ \\

    \theutterance \stepcounter{utterance}  
    & & & \multicolumn{2}{p{0.3\linewidth}}{
        \cellcolor[rgb]{0.9,0.9,0.9}{
            \makecell[{{p{\linewidth}}}]{
                \texttt{\tiny{[GM$|$GM]}}
                \texttt{attempts = 2} \\
\texttt{target\_word = djeca} \\
\texttt{guess = vlada} \\
\texttt{guess\_feedback = v<red> l<red> a<red> d<yellow> a<green>} \\
            }
        }
    }
    & & \\ \\

    \theutterance \stepcounter{utterance}  
    & & & \multicolumn{4}{p{0.6\linewidth}}{
        \cellcolor[rgb]{0.9,0.9,0.9}{
            \makecell[{{p{\linewidth}}}]{
                \texttt{\tiny{[P1$\langle$GM]}}
                \texttt{povratna\_informacija\_pogađanja: v<crvena> l<crvena> a<crvena> d<žuta> a<zelena>} \\
\\ 
\texttt{Daj svoj odgovor samo u ovom formatu.} \\
\texttt{objašnjenje: detalji} \\
\texttt{pogađanje: riječ} \\
            }
        }
    }
     \\ \\

    \theutterance \stepcounter{utterance}  
    & \multicolumn{4}{p{0.6\linewidth}}{
        \cellcolor[rgb]{0.561,0.737,0.561}{
            \makecell[{{p{\linewidth}}}]{
                \texttt{\tiny{[P1$\rangle$GM]}}
                \texttt{objašnjenje: a je fiksna na poziciji 5 (zelena); d je u riječi ali ne na poziciji 3 ni 4, pa mora biti na 1 ili 2; isključena su slova r, i, o, v, l. Biram riječ koja počinje s d i završava na a.} \\
\texttt{pogađanje: desna} \\
            }
        }
    }
    & & \\ \\

    \theutterance \stepcounter{utterance}  
    & & & \multicolumn{2}{p{0.3\linewidth}}{
        \cellcolor[rgb]{0.9,0.9,0.9}{
            \makecell[{{p{\linewidth}}}]{
                \texttt{\tiny{[GM$|$GM]}}
                \texttt{Error: The guessed word is not a valid word for this game.} \\
            }
        }
    }
    & & \\ \\

    \theutterance \stepcounter{utterance}  
    & & & \multicolumn{4}{p{0.6\linewidth}}{
        \cellcolor[rgb]{0.9,0.9,0.9}{
            \makecell[{{p{\linewidth}}}]{
                \texttt{\tiny{[P1$\langle$GM]}}
                \texttt{Tvoje pogađanje nije valjana riječ za ovu igru. Pokušaj ponovno.} \\
\\ 
\texttt{Daj svoj odgovor samo u ovom formatu.} \\
\texttt{objašnjenje: detalji} \\
\texttt{pogađanje: riječ} \\
            }
        }
    }
     \\ \\

    \theutterance \stepcounter{utterance}  
    & \multicolumn{4}{p{0.6\linewidth}}{
        \cellcolor[rgb]{0.561,0.737,0.561}{
            \makecell[{{p{\linewidth}}}]{
                \texttt{\tiny{[P1$\rangle$GM]}}
                \texttt{objašnjenje: trebam riječ koja završava na a (zelena) i sadrži d na poziciji 1 ili 2, bez slova r, i, o, v, l. Biram "djeca".} \\
\texttt{pogađanje: djeca} \\
            }
        }
    }
    & & \\ \\

    \theutterance \stepcounter{utterance}  
    & & & \multicolumn{2}{p{0.3\linewidth}}{
        \cellcolor[rgb]{0.9,0.9,0.9}{
            \makecell[{{p{\linewidth}}}]{
                \texttt{\tiny{[GM$|$GM]}}
                \texttt{attempts = 3} \\
\texttt{target\_word = djeca} \\
\texttt{guess = djeca} \\
\texttt{guess\_feedback = d<green> j<green> e<green> c<green> a<green>} \\
            }
        }
    }
    & & \\ \\

    \theutterance \stepcounter{utterance}  
    & & & \multicolumn{2}{p{0.3\linewidth}}{
        \cellcolor[rgb]{0.9,0.9,0.9}{
            \makecell[{{p{\linewidth}}}]{
                \texttt{\tiny{[GM$|$GM]}}
                \texttt{game\_result = WIN} \\
            }
        }
    }
    & & \\ \\

\end{supertabular}
}

\end{document}
